\section*{Теоретическая справка}

\subsection*{1. Предел функции в точке}
Окрестностью точки называется любой промежуток, содержащий данную точку, причём данная точка обязательно является внутренней точкой выбранного промежутка.

\begin{exercise}
\textbf{Определение предела.} Число $A$ называется пределом функции $y=f(x)$ при $x\to a$, если в окрестности точки $a$ значение функции приближается к числу $A$. Обозначение:
\[
  \lim\limits_{x\to a} f(x)=A.
\]
\textbf{Односторонние пределы.} Предел справа записывают $\lim\limits_{x\to a+0}f(x)=A$.
Предел слева записывают $\lim\limits_{x\to a-0}f(x)=A$.

Если функция имеет в некоторой точке равные пределы справа и слева, то существует и её предел в этой точке:
\[
\lim\limits_{x\to a+0}f(x)=\lim\limits_{x\to a-0}f(x)=A
\quad\Longrightarrow\quad
\lim\limits_{x\to a}f(x)=A.
\]
\end{exercise}

\subsection*{2. Правила вычисления пределов}
Если функция определена в некоторой точке (непрерывна в ней), то её предел в этой точке равен её значению: $\lim\limits_{x\to a}f(x)=f(a)$.

Если пределы $f(x)$ и $g(x)$ при $x\to a$ существуют, то:
\begin{enumerate}
  \item $\lim\limits_{x\to a}c=c$, где $c$ --- константа;
  \item $\lim\limits_{x\to a}(f(x)\pm g(x))=\lim\limits_{x\to a}f(x)\pm\lim\limits_{x\to a}g(x)$;
  \item $\lim\limits_{x\to a}(f(x)g(x))=\lim\limits_{x\to a}f(x)\cdot\lim\limits_{x\to a}g(x)$;
  \item $\lim\limits_{x\to a}\dfrac{f(x)}{g(x)}=\dfrac{\lim\limits_{x\to a}f(x)}{\lim\limits_{x\to a}g(x)}$, если $\lim\limits_{x\to a}g(x)\ne0$.
\end{enumerate}

\subsection*{3. Предел на бесконечности и неопределённости}
При $x\to\infty$ (как к $+\infty$, так и к $-\infty$) значения функции могут стремиться к конкретному числу; например, $\lim\limits_{x\to\infty}\frac1x=0$.

Часто встречаются неопределённости
\[
  \left[\frac\infty\infty\right],\quad \left[\frac00\right],\quad [0\cdot\infty],\quad [\infty-\infty],\quad [0^0],\quad [\infty^0],\quad [1^\infty].
\]
Для их раскрытия используют алгебраические преобразования (разложение на множители, домножение на сопряжённое) и замечательные пределы.

\subsection*{4. Замечательные пределы}
\begin{exercise}
\textbf{Первый замечательный предел} применяется при неопределённости $\left[\frac00\right]$:
\[
  \lim\limits_{x\to0}\frac{\sin x}{x}=1.
\]
Следствия: $\lim\limits_{x\to0}\frac{\sin x}{x}=1$, $\lim\limits_{x\to0}\frac{\tg x}{x}=1$, $\lim\limits_{x\to0}\frac{\arcsin x}{x}=1$, $\lim\limits_{x\to0}\frac{\arctg x}{x}=1$.
\end{exercise}

\begin{exercise}
\textbf{Второй замечательный предел} применяется при неопределённости $[1^\infty]$:
\[
  \lim\limits_{x\to\infty}\left(1+\frac1x\right)^x=e,
  \qquad
  \lim\limits_{x\to0}(1+x)^{1/x}=e.
\]
\end{exercise}
